\documentclass[12pt]{article}

\usepackage{amsmath,amssymb,amsthm}
\usepackage{geometry}
\usepackage{hyperref}
\usepackage{setspace}

\geometry{margin=1in}
\onehalfspacing

\title{Branch Formulas for the Collatz Map\\
\large A First-Principles Derivation via Steiner Circuits}
\author{Jon Seymour\\ \texttt{jon@wildducktheories.com}\\ Sydney}
\date{March 1, 2026}

\begin{document}
\maketitle

\begin{abstract}
We present a complete derivation of the branch formulas \(A(n,x)\), \(B(n,x)\),
and their common endpoint \(C(n,x)\) for the Collatz map. These formulas classify
odd integers according to initial parity patterns and arise naturally from the
closure conditions of canonical Steiner circuits. The derivation is fully
first-principles and independent of any global convergence claim.
\end{abstract}

\section{Introduction}

The Collatz map is defined by
\[
T(n) =
\begin{cases}
n/2, & n \text{ even},\\[2mm]
(3n+1)/2, & n \text{ odd}.
\end{cases}
\]

A fruitful approach to analyzing Collatz trajectories is to group integers
according to their parity patterns rather than numerical magnitude. Odd integers
can be organized into \emph{branches} determined by finite parity sequences that
govern the initial segment of their trajectory.

The branch framework studied here—together with the explicit formulas
\(A(n,x), B(n,x), C(n,x)\)—is inspired by the structural approach introduced
in Neel Alkoraishi’s preprint
\emph{A Bit-Length and Branch-Based Proof of the Collatz Conjecture}, Version~10.\footnote{
Neel Alkoraishi, ``A Bit-Length and Branch-Based Proof of the Collatz Conjecture,''
Version~10, Zenodo (2026),
\url{https://doi.org/10.5281/zenodo.18736864}.}
This work isolates the algebraic content of the branch framework without
adopting any global convergence claims.

\section{Parity Vectors and Steiner Circuits}

A \emph{parity vector} (or \emph{Steiner circuit})\cite{Steiner1977} of length \(m\) is a finite
sequence
\[
\sigma = (\sigma_0, \dots, \sigma_{m-1}), \quad \sigma_i \in \{E,O\},
\]
where \(E\) denotes the map \(x \mapsto x/2\) and \(O\) denotes \(x \mapsto (3x+1)/2\).

Composing the Collatz map along a circuit yields an affine map
\[
x \longmapsto \frac{3^k x + c}{2^m},
\]
where \(k\) is the number of odd steps and \(c\) is determined by the circuit.
Integrality imposes congruence conditions on \(x\), which determine the
admissible odd integers in the branch.

\section{The Canonical Circuit \((OE)^n\)}

We focus on the canonical circuit
\[
(OE)^n = OEOE\cdots OE,
\]
with \(n\) repetitions. This circuit has length \(2n\) and contains exactly
\(n\) odd steps. Composing the Collatz map along this circuit gives
\[
T^{(OE)^n}(N) = \frac{3^n N + K_n}{2^n}, \qquad 
K_n := \sum_{j=0}^{n-1} 3^j 2^{\,n-1-j}.
\]

The integrality condition requires
\[
3^n N + K_n \equiv 0 \pmod{2^n}.
\]

---

\section{Derivation of the Branch Formulas}

\subsection*{Step 1: Integrality condition}

For \(T^{(OE)^n}(N)\) to be an integer, we must satisfy
\[
3^n N + K_n \equiv 0 \pmod{2^n}.
\]

\subsection*{Step 2: Evaluate \(K_n\) modulo \(2^{n+1}\)}

Reducing \(K_n\) modulo \(2^{n+1}\) gives
\[
K_n \equiv
\begin{cases}
2^{n-1}, & n \text{ even},\\[1mm]
3 \cdot 2^{n-1}, & n \text{ odd},
\end{cases}
\pmod{2^{n+1}}.
\]

---

\subsection*{Step 3: Lift solutions modulo \(2^n\) to \(2^{n+1}\)}

Let \(N_0\) be a solution modulo \(2^n\). Then all solutions modulo \(2^{n+1}\)
are
\[
N = N_0 + t \cdot 2^n, \quad t = 0 \text{ or } 1.
\]

This standard lifting procedure shows there are exactly two odd residues modulo
\(2^{n+1}\), corresponding to the two branches \(A\) and \(B\).

---

\subsection*{Step 4: Solve explicitly for residues}

\textbf{Even \(n\)}:
\[
N \equiv 2^{n-1}-1 \quad \text{or} \quad 2^n-1 \pmod{2^{n+1}}.
\]

\textbf{Odd \(n\)}:
\[
N \equiv 3\cdot 2^{n-1}-1 \quad \text{or} \quad 3\cdot 2^n-1 \pmod{2^{n+1}}.
\]

These are the “base” values for the two branches.

---

\subsection*{Step 5: Introduce the enumeration parameter \(x\)}

All integers in the same residue class modulo \(2^{n+1}\) can be written as
\[
N = N_{\text{base}} + k \cdot 2^{\,n+1}, \quad k \ge 0.
\]

- \(N_{\text{base}}\) is one of the two residues from Step 4.  
- \(k\) enumerates all integers in the residue class.  
- Define \(x := k\), which arises naturally from modular arithmetic.

Thus \(x\) is not arbitrary; it indexes all odd integers in the branch.

---

\subsection*{Step 6: Final branch formulas}

\[
A(n,x) =
\begin{cases}
3\cdot 2^{n-1} + 2^{n+1}x - 1, & n \text{ odd},\\[1mm]
2^{n-1} + 2^{n+1}x - 1, & n \text{ even},
\end{cases}
\]

\[
B(n,x) =
\begin{cases}
3\cdot 2^{n} + 2^{n+2}x - 1, & n \text{ odd},\\[1mm]
2^{n} + 2^{n+2}x - 1, & n \text{ even}.
\end{cases}
\]

These formulas enumerate all odd integers beginning with the canonical
\((OE)^n\) circuit.

---

\section{Branch Endpoints}

Substituting \(A(n,x)\) or \(B(n,x)\) into the affine map yields the common
endpoint
\[
C(n,x) = 2 \cdot 3^n x + 4 \sum_{i=0}^{\lfloor (n-1)/2 \rfloor} 9^i.
\]

Properties:

\begin{itemize}
\item \(C(n,x)\) is always even;
\item \(C(n,x)\) is divisible by \(4\);
\item both \(A(n,x)\) and \(B(n,x)\) reach \(C(n,x)\) after exactly \(n\)
odd steps.
\end{itemize}

This endpoint arises directly from the combinatorial and arithmetic structure
of the circuit, independent of any global convergence argument.

---

\section{Scope and Non-Claims}

The derivation above is entirely algebraic:

- No claim is made regarding global convergence of the Collatz map.
- No layer induction or global ranking invariant is assumed.
- The formulas \(A(n,x), B(n,x), C(n,x)\) are mathematically meaningful on
their own.

\section{Conclusion}

We have derived, from first principles, the explicit branch formulas
\(A(n,x), B(n,x)\) and the associated endpoint \(C(n,x)\), using canonical
Steiner circuits and modular arithmetic. This fully captures the structural
content of the branch framework and clarifies the role of the enumeration
parameter \(x\) in indexing all integers in a given branch.

\appendix

\section{Use of AI}

This paper was generated with the assistance of artificial intelligence tools.
The initial draft of the mathematical content, derivations, and structure were
produced using OpenAI's ChatGPT (GPT-4) in a conversation available at:

\begin{quote}
\url{https://chatgpt.com/c/69a3a8d5-0d90-8322-849f-4c43cf6ecc40}
\end{quote}

Subsequent editing, including the addition of author metadata, location,
citation of Steiner's 1977 paper, and this appendix, was performed using
Anthropic's Claude (Claude Sonnet 4.5) via Claude Code on March 1, 2026.

\subsection*{Original Prompts}

The paper was generated through an iterative dialogue with ChatGPT. The
key prompts in chronological order were:

\begin{enumerate}
\item Initial context-setting prompt establishing the mathematical framework,
including definitions of the Collatz map, branch formulae \(A(n,x)\) and
\(B(n,x)\), endpoint formula \(C(n,x)\), binary decomposition insights, and
agreement to focus on the branch framework independently of convergence claims.

\item ``What I would like to do is highlight the relationship between the
\(A,B,C\) and Steiner circuits. I'd also like to provide a complete derivation
for \(A,B,C\) from first principles.''

\item ``Please ensure that you cite the original paper
(\url{https://zenodo.org/records/18736864}) (specifically Version 10) as the
original inspiration for this write up.''

\item ``Can you explain the 3 different styles?'' [referring to citation styles
APA, MLA, and Chicago]

\item ``Ok, let's use Chicago style. Please generate the full LaTeX source as
discussed.''

\item ``I am not convinced by the Step 5 here - I think you need to show how
\(x\) is derived from what comes before.'' [requesting clarification of the
derivation of the enumeration parameter]
\end{enumerate}

\subsection*{AI Contributions}

\textbf{ChatGPT} (initial generation):
\begin{itemize}
\item Generated the complete mathematical derivation and proof structure
\item Formulated the abstract, introduction, and conclusion
\item Produced the LaTeX formatting and document organization
\item Developed the step-by-step presentation of the branch formula derivation
\end{itemize}

\textbf{Claude Code} (subsequent editing):
\begin{itemize}
\item Added author information (Jon Seymour, \texttt{jon@wildducktheories.com}, Sydney)
\item Set document date to March 1, 2026
\item Researched and added the citation for Steiner's 1977 paper
\item Created bibliography section
\item Configured Makefile build system for PDF generation
\item Added this AI usage disclosure appendix
\end{itemize}

All mathematical content, reasoning, and exposition were AI-generated.
Human input consisted of high-level direction, verification of mathematical
correctness, and editorial decisions.

\begin{thebibliography}{9}

\bibitem{Steiner1977}
R.~P.~Steiner,
\emph{A theorem on the syracuse problem},
Proceedings of the 7th Manitoba Conference on Numerical Mathematics,
Utilitas Mathematica Publishing, Winnipeg, MB, Canada,
pp.~553--559, 1977.

\end{thebibliography}

\end{document}